\section{Lorentz Safety}

Relativity is the hardest test a discrete theory of space has to pass. The deep fact of our universe is
Lorentz invariance. The laws of physics look the same no matter how fast one moves, with no special rest
frame and no preferred direction in space. This is the make-or-break filter for any theory that builds
spacetime from a discrete web of relations, and most such theories fail it in a predictable way.

This section
explains why a flat regular lattice cannot carry relativity, why a hyperbolic and effectively aperiodic
substrate can, how the substrate passes a genuine boost test, how the observational bounds come out, and how
the result reappears as an emergent symmetry of the dynamics.

\subsection{The causal-set lesson}

Causal-set theory taught the field a hard lesson. Build spacetime from a tidy flat lattice, integer points in
neat rows and columns with causal links between them, and the construction is fatal. A flat lattice has
preferred directions. Motion along a row is geometrically different from motion along a diagonal, and because
the space is flat those directions are globally defined. North here is the same as north a mile away.

The
grain is real and detectable everywhere at once, which gives the lattice a preferred rest frame, exactly what
relativity forbids. The damage is concrete and measurable. On a lattice the speed of a signal depends on
which way it travels, and the dependence grows with energy, since high-energy signals feel the fine grain
most. This is the energy-dependent, directional Lorentz violation measured directly in our work (P27). A flat
substrate is not subtly wrong. It announces its preferred frame.

\subsection{Why curvature buys Lorentz safety}

The escape is not what intuition suggests. One would expect that hiding the grain requires a messy random
space and that an exact ordered crystal would be the worst case. The opposite is true. The rescue is
curvature, and an exact ordered hyperbolic crystal is Lorentz-safe. In flat space a direction can be carried
unchanged across the whole space, so a flat grid's preferred directions are globally real. In curved space it
cannot. Curvature scrambles directions under parallel transport. Move across a hyperbolic space and the links
that pointed one way locally now point every which way relative to a distant patch. From any local spot the
crystal's links fan out in all directions, statistically indistinguishable from a random foam with no
preferred direction. The order is still there. It just does not line up across the space, so no resident can
use it to define a global frame.

The experiments confirm this across a wide family of substrates. A fully deterministic substrate, the
golden-angle sunflower, is as Lorentz-safe as a random sprinkle, with directional anisotropy in the same low
band and no random number generator used anywhere (P40, see also P39). Determinism is fine. What matters is
that the sequence is effectively aperiodic, never repeating a global direction.

A sweep across random
sprinkles, the sunflower, a low-discrepancy sequence, and the regular $\{7,3\}$ and $\{5,4\}$ hyperbolic
tilings finds every hyperbolic substrate Lorentz-safe with anisotropy below threshold, while the flat lattice
control is the only one that fails (P40). Regularity does not break Lorentz invariance once the space is
curved. The Margenstern families of hyperbolic tilings, both the $\{p,4\}$ and $\{p,3\}$ series, are all
Lorentz-safe (P41). Lifting into three dimensions, the $\{5,3,4\}$ dodecagrid at the heart of the theory is
Lorentz-safe while a flat cubic lattice is not (P45).

A natural worry is that the dodecagrid is a lucky
special case hand-picked to pass. Feeding $\{7,3\}$, $\{5,4\}$, $\{8,3\}$, $\{6,4\}$, and the 3D $\{5,3,4\}$
through a single Coxeter generator, changing only the Schlafli symbol, confirms every one is Lorentz-safe
(P47). These tilings are one machine on different settings, so Lorentz safety is a property of the whole
reflection-group family, not a fluke of one member.

\subsection{Boost-covariance on the substrate}

Isotropy of directions is only the rotational half of Lorentz invariance. The other half is boost safety,
invariance under changes of velocity, and it must be tested on its own. The boost parameter of a timelike
causal link is its rapidity, the hyperbolic angle of the link in spacetime, and a real boost shifts every
rapidity by the same amount because rapidities add. So a structure with no preferred frame has a rapidity
distribution that is flat and broad, while a structure with a preferred frame has one peaked at rapidity zero
where its timelike links pile up along a built-in time axis.

A Poisson sprinkle of $1+1$ Minkowski space gives
a flat distribution with normalized entropy above $0.9$, while the matched causal lattice gives a sharply
peaked one with a spread several times smaller, betraying its rest frame (P84). The same experiment then
performs an actual boost, applying a Lorentz transformation of rapidity one to every coordinate, leaving the
causal order untouched, and recomputes the rapidities in the new frame.

The distribution shifts by exactly
the boost, as it should, but its shape is unchanged. After removing the expected shift, the flatness of the
boosted distribution matches the original to within a small tolerance. The sprinkle is boost-covariant. Its
statistical description is the same in the boosted frame as in the rest frame, which is precisely what no
preferred frame means. This is a genuine boost test, not a rotation proxy, and the substrate passes it while
the lattice does not (P84).

\subsection{Observational bounds}

Lorentz safety is also an observational check, and the observation is real. Light from distant gamma-ray
bursts travels billions of years before reaching us. If spacetime had a grain, high-energy and low-energy
photons would travel at slightly different speeds and would drift apart over that immense baseline. They
arrive neck and neck, ruling out any first-order, energy-dependent speed variation to extraordinary
precision.

The model's predicted linear-order Lorentz-violation coefficient goes to zero, comfortably inside
the gamma-ray-burst bound, because the curved crystal predicts the grain simply is not there, while the same
calculation on a flat lattice gives a large coefficient that the data excludes (P82). The residual causal-set
swerve, a diffusion of particle momenta, also vanishes as the discreteness scale shrinks, so there is no
leftover Lorentz-violating jitter (P82).

This cuts both ways. The model predicts no first-order grain in the
speed of light, and if anyone ever measured a first-order direction- or energy-dependence of light's speed,
the curved-crystal picture would be falsified. The current data is on its side.

\subsection{The link to the emergent dynamics}

The substrate-level results show the foundation has no built-in frame. The full theory then has to grow
relativity out of its dynamics, not just its geometry, and three modern results complete the chain.

Leading-order rotational isotropy emerges from the dynamics, with the one-step diffusion tensor on the $\{5,3,4\}$
having equal eigenvalues, so a walker spreads the same in every direction. This isotropy is inherited from
the icosahedral symmetry of the dodecagrid's cell, the geometric reason the lattice has no preferred spatial
axis even though it is built from identical cells (P124). The deterministic reversible wave on the
$\{5,3,4\}$ then has an isotropic speed, with anisotropy across the twelve directions consistent with zero, an
emergent frame-independent light speed built from a real dynamical rule rather than from substrate statistics
alone (P150).

The boost rung follows on the same dynamics. The wave's dispersion is real, linear, and
massless, $\omega = |k|$, a positive-norm massless mode with an exact lightcone that is boost-invariant, with
massive modes boost-invariant in the infrared window below the cutoff (P154). Together with the rotations,
this gives the full emergent Lorentz group acting on a genuine relativistic field.

The arc is complete. The substrate has no preferred frame (P27, P39, P40, P41, P45, P47, P84). The
observations that kill a lattice are passed (P82). And relativity reappears as an emergent symmetry of the
actual dynamics on the dodecagrid (P124, P150, P154).

The choice of a negatively curved substrate was made
for several reasons at once. Curvature gives room to grow, it carries the right symmetries, and it removes
the preferred frame that flat discreteness cannot escape. Room, symmetry, and Lorentz invariance are not
three separate gifts. They fall out of a single decision. Curvature, not disorder, buys Lorentz invariance,
and the $\{5,3,4\}$ dodecagrid is exactly the curved, aperiodic substrate that relativity needs.
